\documentclass{jsarticle}
\usepackage[dvipdfmx,final]{graphicx}
\usepackage{amsmath}
\usepackage{amssymb}
\begin{document}
\title{時間計量の微分変数についての大域的微小量の表し方}
\author{Masaaki Yamaguchi}
\date{}
\maketitle
\newcommand{\timebow}{%
	\begin{picture}(5,5)
		\put(0,0){$\text{\rotatebox{90}{$\bowtie$}}$}
	\end{picture} %
	}
   \newcommand{\variint}{%
	\begin{picture}(15,30)
		\put(0,0){
			$ \iint $}
		\put(0,5){\line(15,0){15}}
	\end{picture} %
}

\newcommand{\variiint}{%
	\begin{picture}(15,15)
		\put(0,0){
			$ \iiint $}
		\put(0,5){\line(15,0){20}}
	\end{picture} %
}



 \newcommand{\alearletter}{%
	 \begin{picture}(10,20)
		 \put(0,0){
			 $ \square $}
		 \put(0,0){\line(1,1){10}}
	 \end{picture} %
	 }

 重力場理論の式は、
  Gravity equation is
  $$ \square = - {16 \pi G \over c^4}T^{\mu\nu} $$
  This equation quated with being logment of formula, and
  this formula divided with universe of number in prime zone,
  therefore, this dicided with varint equation is monotonicity of
  being composited with Weil's theorem united for Gamma function.
  この式は、対数を宇宙における数により求める素数分布論として、この
  大域的積分分断多様体がガンマ関数をヴェイユ予想を根幹とする単体量
  として決まることに起因する商代数として導かれる。
  $$ \scalebox{1}[2]{\alearletter} = 
  {{{{8 \pi G} \over {c^4}}T^{\mu\nu}}/ \log x} $$
  $$ {}^{t}\scalebox{1}[2]{\variiint}\mathrm{cohom D_{\chi}[I_m]}$$
  $$ = \oint {(px^n + qx + r)^{\nabla l}} $$
  $$ {d \over d l}L(x,y) = 2 \int ||\sin 2x||^{2}d \tau $$
  $$ {d \over d \gamma}\Gamma $$
  この関数は大域的微分多様体としてのアカシックレコードの合流地点として、
  タプルスペースを形成している。
  This function esterminate with acasic record of global differential
  manifolds.
  $$ = [i \pi (\chi, x), f(x)] $$
  それにより、この多様体は基本群をアカシックレコードの相対性としての
  存在論の実存主義から統合される多様体自身としてのタプルスペースの
  池になっている。
  And this manifold from fundemental group esterminate with
  also this manifold estimate relativity of acasic record.
  $$ {d \over d \gamma}\Gamma = \int C dx_m = \int(\int {1 \over x^s}dx
  - \log x)\mathrm{dvol} $$
  また、このアカシックレコードはオイラーの定数のラムダドライバー
  にもなっている。
  More also this record tupled with lake of Euler product.
  $$ \bigoplus({i \hbar^{\nabla}})^{\oplus L} $$
  $$ x^{{1 \over 2} + iy} = e^{x \log x} $$
  それゆえに、この定数はゼータ関数と微分幾何の量子化を因数にもつ
  素因子分解の式にもなっている。
  Therefore, this product also constructed with differential geometry
  of quantum level equation and zeta function.
  $$ = C $$ 
  そして、この関数はオイラーの定数から広中平祐定理による4重帰納法の
  オイラーの公式からの多様体積分へとのサーストン空間のスペクトラム関数
  ともなっている。
  And this function of Euler product respectrum of focus with
  Heisuke Hironaka manifold in four assembled of integral Euler equation.. 

  $$ C =  b_0 + \cfrac{c_1}{b_1 +
	  \cfrac{c_2}{b_2 +
	  \cfrac{c_3}{b_3 +
	  \cfrac{c_4}{b_4 + \cdots}}}} $$
  この方程式は指数による連分数としての役割も担っている。  
  This equation demanded with continued number of step function.
  $$ x^{{1 \over 2} + iy} = e^{x \log x} $$
  $$ (2.71828)^{2.828} = a^{a + b^{b + c^{c + d^{d \cdots }}}} $$
  $$ = \int e^f \cdot x^{1 - t}dx $$
  This represent is Gamma function in Euler product.
  Therefore this product is zeta function of global differential equation.

  $$ {d \over d f}F(v_{ij},h) = \int {e^{-f}\mathrm{dV}}
  [- \Delta v + \nabla_{i}\nabla_{j}
  v_{ij} - R_{ij}v_{ij} - v_{ij}\nabla_{i}\nabla_{j} + 2 <\nabla f,\nabla h>
  + (R + \nabla f^2)({v \over 2} - h)]$$
  $$ = {d \over d f}F = {2 {\int (R + \nabla_{i}\nabla_{j} f)^2} \over 
  { - (R + \Delta f)}}\mathrm{dm} $$
  これらの方程式は8種類の微分幾何の次元多様体として、そして、これらの
  多様体は曲平面による双対性をも生成している。 
  そして、このガウスの曲平面は、大域的微分多様体と微分幾何の量子化から
  素因数を形成してもいる。
  These equation are eight differential geometry of dimension calyement.
  And these calyement equation excluded into pair of dimension surface.
  This surface of Gauss function are global differential manifold,
  and differential geometry of quantum level.
  $$ F \ge {d \over d f}\int \int{1 \over (x \log x)^2}dx_m 
  + {d \over d f}\int \int{1 \over (y \log y)^{1 \over 2}}dy_m $$

  微分幾何の量子化はオイラーの定数とガンマ関数が指数による連分数
  としての不変性として素因数を形成していて、このガウスの曲平面
  による量子力学における重力場理論は、ダランベルシアンの切断多様体
  がこの大域的切断多様体を付加してもいる。
  Differential geometry of quantum level constructed with Euler product
  and Gamma function being discatastrophed from continued fraction style.
  Gravity equation lend with varint of monotonicity of level expresented
  from gravity of letter varient formula.
  これらの方程式は基本群と大域的微分多様体をエスコートしていもいて、
  ヴェイユ予想がこのダランベルシアンの切断方程式たちから
  輸送のポートにもなっている。ベータ関数とガンマ関数がこれらのフォームラ
  の方程式を放出してもいて、結果、これらの方程式は広中平祐定理の複素
  多様体とグリーシャ教授によるペレルマン多様体からサポートされてもいる。
  この2名の教授は、一つは抽象理論をもう一方は具象理論を説明としている。
  These equation escourted into Global differential manifold and
  fundemental equation.
  Weil's theorem is imported from this equation in gravity of letters.
  Beta function and Gamma function are excluded with these formulas.
  These equation comontius from Heisuke Hironaga of complex manifold
  and Gresha professor of Perelman manifold.
  These two professos are one of abstract theorem and the other of
  visual manifold theorem.


ガンマ関数についての大域的部分積分多様体は、単体量の時間計量に、
$ \timebow_m $を微分変数の微小量に使う

$$ {T^{\mu\nu}}^{T^{\mu\nu}} = \int T^{\mu\nu} \timebow_m $$
   $$ {T^{\mu\nu}}^{T^{\mu\nu}} = \int T^{\mu\nu} \rotatebox{90}{$\bowtie$}_m $$
   ガンマ関数についての大域的微分多様体は
   $$   {T^{\mu\nu}}^{{T^{\mu\nu}}^{'}}=\int T^{\mu\nu}dx_m  $$
      $$ {\Gamma^{\gamma}}^{'} = {d \over {d \gamma}}\Gamma $$
      大域的積分多様体は
   $$ \Gamma^{\gamma} = \int \Gamma dx_m $$
   このxのx乗の大域的積分多様体は　
   $$ \rotatebox{90}{$\bowtie$} = \int x^x dx_m $$

一般相対性理論の多様体積分が$ {T^{\mu\nu}}^{T^{\mu\nu}}$と表せられて
  $$ {T^{\mu\nu}}^{T^{\mu\nu}} = \int T^{\mu\nu}
   \timebow_m $$
   $$   {T^{\mu\nu}}^{{T^{\mu\nu}}^{'}}=\int T^{\mu\nu}dx_m  $$
      $$ {\Gamma^{\gamma}}^{'} = {d \over {d \gamma}}\Gamma $$
   $$ \Gamma^{\gamma} = \int \Gamma dx_m $$
   $$ \timebow = \int x^x dx_m $$
   xのxの冪乗の積分についての微小量として、使われる

 $$ T = \int \Gamma(\gamma)^{'}dx_m $$ 
 Jones多項式についてのフェルマーの定理と同型であり、　
 $$ {d \over {d \mathcal{T}}} T^{\mu\nu} = {d \over {d \mathcal{T}}}
 \int T dx_m $$
 ここで、ガンマ関数における大域的部分積分多様体の特異点方程式は
 $$ ||ds^2|| = \int [T]dx_m $$
 と特異点の方程式は、ガウス記号を使う

 これらをまとめると、次の式たちにまとめられることができる
 $$ = \int {\Gamma(\gamma)^{'} \over {h\nu}}dx_m $$
 $$ = \int{\Gamma(\gamma)^{'} \over {x \log x}}dx_m $$
 $$ = d\timebow_m  $$
 $$ {d \over {d \mathcal{T}}}\int T^{\mu\nu}dx_m $$
 $$ = {d \over {d \mathcal{T}}}T^{\mu\nu }$$
 波動関数の波長は、光量子仮設で対数を同型として
 $$ \lambda = h \nu $$

 これらから、対数が
 $$\sqrt{g}=1={1 \over {\log x}} $$
 と、ゼータ関数と関係していることが言える

 $$ y = x - a, g = \log x = \log a $$
 差関数は
 $$ x = a $$
 この数値たちは、対数関数で
 $$ a = {a \over {\log a}} = 1 = {1 \over {\log a}} $$
 $$ = \sqrt{g} = 1 $$
 と、ゼータ関数の基盤になっていて
 ゼータ関数を対数と指数関数とで表せると
 $$ x^{{1 \over 2} + iy} = x^{{d \over {df}} \int\int{1 \over (y \log y)^
 {1 \over 2}}dx_y} $$
 ベータ関数として
 $$ {{L^{n+1}} \over {n+1}} = \int x^{1 - t}t^{x -1}dt $$
 $$ \beta(p,q) $$
 ゼータ関数は、
 $$ {d \over {df}}\int\int{1 \over {(y \log y)^{1 \over 2}}}dy_m 
 \ge {1 \over 2} $$
 $$ {d \over {df}}F = F^{f^{'}} $$
 $$ = e^{x \log x} $$
 フェルマーの定理としての、大域的積分多様体は
 $$ \int F dx_m = e^{2x \log x} $$
 $$ = F^f = e^{x \log x} + e^{-x \log x} \ge e^{2 x \log x} $$
 $$ x^{x^{x \log x}} = x^{x^x} $$
 全数値は、直線体であり、
 $$ 2^2 \cong 3^3 \cong 5^5 \cong 7^7 \cong 11^{11} \cong 13^{13} \cong
 17^{17} \cong 23^{23} $$
 冪乗が対数として
 $$ x^x = a, x \log x = \log a, x = {\log a \over {\log x}} $$
 $$ x = \log {a \over x} $$
 $$ \sqrt{g} = 1, \sqrt{a} = x, x = y^0, x^x = e^{x \log x} $$
 と、ゼータ関数になり
 $$ a = {a \over {\log a}} = 1 = {1 \over {\log a}} $$
 $$ x^x = a, e^x = a^L = {1 \over x}, x \to \infty \to 0 $$
 $$ e^{-f}dV = {f \over {\log x}} $$
 これらの式たちは、ゼータ関数への導く式たちである
 
 $$ \delta \int (L - {1 \over {2\kappa}}R)\sqrt{g} d^4 x = 0 $$
 $$ R^{\mu\nu} = - \kappa(T^{\mu\nu} - {1 \over 2}g^{\mu\nu}T) $$ 
 $$ R^{\mu\nu} - {1 \over 2}g^{\mu\nu}R = - \kappa T^{\mu\nu} $$
 $$ (R^{\mu\nu} - {1 \over 2}g^{\mu\nu}R)_{i\nu} = 0$$
 $$ {\partial \over {\partial x^{\alpha}}}(g^{\alpha \beta}{\partial
 \log \sqrt{g} \over {\partial x^{\beta}}}) = - \kappa T $$
 
  これらの式は、次の式たちへと行く
  $$ \square \Psi = - {16 \pi G \over c^4}T^{\mu\nu} $$
  $$ = {\square \Psi \over {\log x}} = {d \over df}\int \int
  {1 \over {(y \log y)^{1 \over 2}}}dy_m $$
  $$ \scalebox{1}[2]{\alearletter} = 
  {{{{8 \pi G} \over {c^4}}T^{\mu\nu}}/ \log x} $$
  $$ {}^{t}\scalebox{1}[2]{\variiint}\mathrm{cohom D_{\chi}[I_m]}$$
  $$ = \oint {(px^n + qx + r)^{\nabla l}} $$
  $$ {d \over d l}L(x,y) = 2 \int ||\sin 2x||^{2}d \tau $$
  $$ {d \over d \gamma}\Gamma $$
 

\end{document}
